% !TEX root = ../presentation.tex
% =============================================================
%  Frame — definition and physical origin of R', L', G', C'
%  Part of ac-transmission; \input from ac-transmission.tex,
%  the opening frame of the chapter.
%
%  Defines the four per-unit-length parameters that Z' and Y' are
%  built from, and names the physical effect behind each.  Figure
%  fig/line-parameters.tex (inline TikZ) is the nominal pi-section
%  written in those parameters; two-port.tex then abstracts it.
%  Values: 500 kV reference line, 3x Bluebird ACSR at 75 C, 60 Hz.
%  G' is the order of magnitude quoted for OHL in
%  power-transmission.tex; it is negligible against omega*C'.
% =============================================================
\providecommand{\acfig}[1]{ac-transmission/fig/#1}

\begin{frame}{The line parameters $R'$, $L'$, $G'$ and $C'$}
  \begin{block}{Per unit length}
    {Four parameters describe the whole line\\
    Two in series, two in shunt --- each a distinct physical effect.}
  \end{block}

  \vspace{0.3em}
  \begin{columns}[c,onlytextwidth]
    \begin{column}{0.44\textwidth}\centering
      \input{\acfig{line-parameters.tex}}\\[1.0em]
      \setlength{\abovedisplayskip}{2pt}\setlength{\belowdisplayskip}{2pt}
      {\small
      \[ \begin{aligned}
           \underline{Z}' &= R'+j\cdot\omega\cdot L' \\
           \underline{Y}' &= G'+j\cdot\omega\cdot C'
         \end{aligned} \]}
    \end{column}
    \begin{column}{0.53\textwidth}
      \renewcommand{\arraystretch}{1.1}
      {\footnotesize
      \begin{tabularx}{\linewidth}{@{}l l X r@{}}
        & \textbf{Parameter} & \textbf{Effect}
          & \multicolumn{1}{c}{\textbf{Value}} \\
        \hline
        $R'$ & Resistance  & Ohmic loss      & 56.9\,m$\Omega$/km \\
        $L'$ & Inductance  & Magnetic field  & 0.891\,mH/km \\
        $G'$ & Conductance & Leakage         & 0.05\,\textmu S/km \\
        $C'$ & Capacitance & Electric field  & 12.7\,nF/km \\
        \hline
      \end{tabularx}}\\[0.6em]
      {\color{atGray}\footnotesize From the same geometry as $X$:
      $L'=\frac{\mu_0}{2\cdot\pi}\cdot\ln\frac{D_m}{r_b}$,
      $C'=\frac{2\cdot\pi\cdot\varepsilon_0}{\ln\left(D_m/r_{b\_c}\right)}$.
      On an OHL $G'\ll\omega\cdot C'=4.79$\,\textmu S/km.}
    \end{column}
  \end{columns}
\end{frame}
